%------------------------------------------------------------------------------%
\section{Área Entre Curvas}
\label{sec:Area_Entre_Curvas}

Ao definirmos integral definida, definimos também o que seria a área sob
uma curva para uma função não negativa. Aqui, usaremos integrais para
encontrar áreas de regiões entre gráficos de duas funções.

Considere que queremos calcular a área da região $S$ entre duas curvas
$y=f(x)$ e $y=g(x)$ e as retas verticais $x=a$ e $x=b$.

\begin{figure}
  \centering
  \includegraphics[width=0.5\linewidth]{area_entre_curvas_01}
  \caption{Gráfico da região.}
  \label{fig:area_entre_curvas_01}
\end{figure}

Vamos dividir a região $S$ em $n$ faixas de mesma largura e então
aproximamos a $k$-ésima faixa por um retângulo com base
$\Delta x$ e altura $f(c_k)- g(c_k)$, conforme mostra a
Figura~\ref{fig:area_entre_curvas_07}.

\begin{figure}
  \centering
  \includegraphics[width=0.5\linewidth]{area_entre_curvas_07}
  \caption{Gráfico da região.}
  \label{fig:area_entre_curvas_07}
\end{figure}

Podemos aproximar a área da região $S$ como soma da área desses $n$
retângulos
\[
  \sum_{k=1}^{n}[f(c_k)-g(c_k)]\Delta
\]

Essa aproximação fica cada vez melhor quando $n \to \infty$ e assim,
definimos a área da região $S$ como sendo
\[
  A= \displaystyle \lim_{k \to \infty}\sum_{k=1}^{n}[f(c_k)-g(c_k)]\delta x
\]

Reconhecemos o limite acima como a integral definida de $f-g$.
Assim, podemos apresentara definição a seguir.

\begin{definicao}{Área entre curvas}{def:area-entre-curvas}
  A área $A$ da região limitada pelas curvas $y=f(x)$ e $y=g(x)$ e
  pelas retas $x=a$, $x=b$, onde $f$ e $g$ são contínuas e
  $f(x) \geq g(x)$ para todo $x \in [a,b]$, é
  \[
    A= \int_a^b [f(x)-g(x)]dx
  \]
\end{definicao}

%------------------------------------------------------------------------------%
